\documentclass[11pt]{article}

\usepackage[margin=1in]{geometry}
\usepackage{amsmath,amssymb,amsthm}
\usepackage{booktabs}
\usepackage{enumitem}
\usepackage[T1]{fontenc}
\usepackage{lmodern}
\usepackage{xurl}
\usepackage[hidelinks]{hyperref}
\hypersetup{
  pdftitle={An Exact Radial Obstruction in a Base-19 Reverse-and-Add Constraint System},
  pdfauthor={Anonymous}
}

% Generated; do not edit.
\newcommand{\BaseValue}{19}
\newcommand{\CandidateEquationCount}{409}
\newcommand{\CoreEquationCount}{408}
\newcommand{\RadialLayerCount}{35}
\newcommand{\TerminalLayer}{33}
\newcommand{\PeakReachableStates}{314}
\newcommand{\MaximumFrontierDigits}{2}
\newcommand{\MaximumFrontierCarries}{12}
\newcommand{\ActiveDigitCount}{409}
\newcommand{\ActiveCarryCount}{410}
\newcommand{\DeletionAuditCount}{409}
\newcommand{\SatDeletionCount}{408}
\newcommand{\UnsatDeletionCount}{1}
\newcommand{\RemovedCoreEquation}{\texttt{phase=0,n=0,column=82}}
\newcommand{\KnownSatRemovedEquation}{\texttt{phase=0,n=0,column=0}}
\newcommand{\CoreReplaySeconds}{34.53}
\newcommand{\CertificateGenerationSeconds}{172.90}
\newcommand{\DeletionAuditSeconds}{32677.27}
\newcommand{\CleanCloneSeconds}{245.46}
\newcommand{\LeanWitnessSeconds}{37.60}
\newcommand{\CandidateCertificateSha}{\texttt{1ce1b3c59bd95264e5338f7fba5bdffa15fb7635917964980057615395cf138c}}
\newcommand{\CoreCertificateSha}{\texttt{7635e7188197021cfc15cbad6d087f6fdbfa7a8bbfe4e4ed9b91a15a89498c67}}
\newcommand{\DeletionAuditSha}{\texttt{0566f6e2025aec28fe0d2f824f1f6b343cf912e9191a9036629a2a1d4ef03325}}
\newcommand{\GenericSoundnessTheorem}{\texttt{Base19Radial.unsat\_of\_closed\_empty\_frontier}}
\newcommand{\ConcreteSatTheorem}{\texttt{Base19Radial.deleteP0C0Sat}}


\newtheorem{theorem}{Theorem}
\newtheorem{proposition}{Proposition}
\newcommand{\Ccandidate}{\ensuremath{\mathcal{C}}}
\newcommand{\Ccore}{\ensuremath{\mathcal{C}^{\prime}}}
\newcommand{\Sat}{\mathsf{SAT}}
\newcommand{\Unsat}{\mathsf{UNSAT}}

\title{An Exact Radial Obstruction in a Base-19 Reverse-and-Add Constraint System}
\author{Anonymous}
\date{August 2026}

\begin{document}
\maketitle

\begin{abstract}
We study one explicitly frozen bounded-integer constraint system extracted
from a six-phase reverse-and-add topology in base~\BaseValue.  Exact radial
frontier enumeration gives an empty terminal frontier for a
\CandidateEquationCount-equation candidate and for a
\CoreEquationCount-equation subset \Ccore.  A complete one-equation-deletion
audit supplies a satisfying assignment after deletion of each equation of
\Ccore, so \Ccore is deletion-minimal unsatisfiable in the computational
certificate model.  A standalone checker reconstructs every local
quantifier-free linear-integer problem, enumerates every output projection,
and compares the complete transition relation with the promoted certificate.
Lean~4 checks the generic frontier soundness theorem and, separately, one
concrete satisfying neighbor.  The two concrete unsatisfiability results are
\emph{not kernel-checked}: they rely on exact Z3 enumeration and do not carry
an LRAT-like solver proof.  The result is not a minimum-cardinality-core
claim, gives no all-base theorem, and does not establish any general Lychrel
claim.
\end{abstract}

\section{Introduction and Scope}

For a base-$b$ word, the reverse-and-add operation adds the represented
integer to the integer represented by the reversed word.  Carries couple
columns that are adjacent in the usual direction, while reversal couples
digits that are far apart.  A natural way to expose the resulting structure
is therefore to process columns radially, from both ends toward the center.

This paper reports a finite result about one frozen system, not about the
unbounded reverse-and-add process.  The topology was motivated by David
Seal's reported period-six self-similar construction in base 17
\cite{brownSeal}.  The OEIS sequence A060382 records conjectural least
palindrome-free starts across bases and explicitly warns that, apart from its
base-two entry, its listed values are unproved \cite{oeisA060382}.  Neither
source is used as a premise of the result below, and the present obstruction
does not settle the base-19 entry of that sequence.

Our contribution is an auditable artifact for the following scoped
statement.  The fixed candidate \Ccandidate{} is unsatisfiable, and deleting
one specified equation produces a deletion-minimal unsatisfiable subset
\Ccore.  ``Deletion-minimal'' means that removing any one remaining equation
makes this particular subset satisfiable; it does not mean that no smaller
unsatisfiable subset exists.  The concrete claims have exact computational
replays, whereas the general soundness implication has a Lean proof.  We keep
these assurance levels separate throughout.

\section{Fixed Constraint System}

Let $b=\BaseValue$.  A digit variable ranges over
$D=\{0,\ldots,b-1\}$, a carry variable over $K=\{0,1\}$, and the first
digit of each phase is constrained to be nonzero.  There are six phases
$p\in\{0,\ldots,5\}$.  Each phase word has three variable blocks.  Phase zero
has block lengths $(26,25,31)$; every other phase has lengths $(29,23,31)$.
For a nonnegative parameter $n$, two constant runs are inserted between the
blocks.  Their digits are $(0,b-1)$ in phases zero and two through five, and
$(b-1,0)$ in phase one; each run has length $n$.  Write the resulting word as
$W_p(n)$.

For $p<5$, phase $p$ relates $W_p(0)$ to $W_{p+1}(0)$.  The last phase
relates $W_5(0)$ to $W_0(1)$, which is the self-similar growth step.  At
column $k$, let $\ell_{p,k}$ and $r_{p,k}$ be the two source atoms selected
from opposite ends, let $t_{p,k}$ be the corresponding target atom, and let
$c_{p,k-1},c_{p,k}\in K$ be the incoming and outgoing carries.  Missing
end atoms and boundary carries are zero.  The column equation is
\begin{equation}\label{eq:column}
  \ell_{p,k}+r_{p,k}+c_{p,k-1}
  =t_{p,k}+b\,c_{p,k}.
\end{equation}

The canonical builder retains columns $0$--$33$ and $48$--$82$ in phase
zero; columns $0$--$33$ and $49$--$82$ in phases one through four; and
columns $0$--$33$ and $49$--$83$ in phase five.  It then omits phase-five
column 82.  This gives the candidate \Ccandidate.  The independent equation
manifest contains exactly \CandidateEquationCount{} equations over
\ActiveDigitCount{} active digit variables and \ActiveCarryCount{} carry
variables.  The retained equations are a subset of the full topology, so
their inconsistency obstructs that topology; no converse encoding claim is
needed.

The equation manifest serializes the base, topology digest, and every atom,
equation, provenance label, and equation order.  Domain and leading-digit conventions are explicit in both checkers and are derived deterministically
from the base and variable names; they are not separate fields in the version
one JSON schema.  A standard-library-only Python builder and a Lean builder
produce canonical equation renderings that are byte-compared during
reproduction, while construction tests check variable counts and model
validity.  These checks guard against statement drift; they do not themselves
prove unsatisfiability.

\section{Radial Decomposition}

Partition the retained equations into ordered radial layers
$L_0,\ldots,L_m$.  At a cut after $L_i$, a variable is \emph{live} precisely
when it occurs in an equation at or before the cut and also in an equation
after the cut.  Let $I_i$ be the canonically sorted live-variable list.  A
frontier state assigns values only to $I_i$.  Variables used only inside one
layer are existentially hidden.

For an input state $s$ on $I_{i-1}$, define
\[
  T_i(s)=\left\{\pi_{I_i}(a):
  a\text{ extends }s\text{ and satisfies every equation in }L_i\right\}.
\]
Starting with $F_{-1}=\{()\}$, exact propagation sets
$F_i=\bigcup_{s\in F_{i-1}}T_i(s)$.  This is variable elimination along a
path decomposition.  Its cost is governed by the live interface rather than
the total number of variables, echoing the classical bounded-width view of
constraint satisfaction \cite{Freuder1990,DechterPearl1989}.

For \Ccore, the promoted certificate contains \RadialLayerCount{} serialized
layers, labeled from $-1$ through \TerminalLayer.  A cut contains at most
\MaximumFrontierDigits{} live digits and \MaximumFrontierCarries{} live
carries.  The largest reachable frontier has \PeakReachableStates{} states,
and the output after layer \TerminalLayer{} is empty.  These figures are read
from the promoted certificate into this manuscript; they are not copied as
independent literals.

\section{Soundness and Certification}

Exact enumeration is stronger than necessary for soundness.  Suppose
$G_{-1},\ldots,G_m$ are sets of interface states such that the empty input is
in $G_{-1}$ and, for every layer,
\begin{equation}\label{eq:closure}
  s\in G_{i-1}\ \wedge\ u\in T_i(s)\quad\Longrightarrow\quad u\in G_i.
\end{equation}

\begin{theorem}[closed empty frontier]\label{thm:soundness}
If every global satisfying assignment restricts to a locally satisfying
assignment in every layer, the frontiers satisfy~\eqref{eq:closure}, and
$G_m=\varnothing$, then the global constraint system is unsatisfiable.
\end{theorem}

\begin{proof}
Assume a global satisfying assignment $a$.  Its restriction to the empty
initial interface is in $G_{-1}$.  Inductively, local satisfaction in $L_i$
puts $\pi_{I_i}(a)$ in $T_i(\pi_{I_{i-1}}(a))$, and closure puts that
restriction in $G_i$.  At the final layer this says
$\pi_{I_m}(a)\in G_m=\varnothing$, a contradiction.
\end{proof}

Lean~4 formalizes this argument as \GenericSoundnessTheorem.  It also checks a
hand-written two-layer contradiction, thereby exercising the theorem with
concrete data.  The reported axiom dependency is only Lean's standard
\texttt{propext} principle \cite{deMouraUllrich2021}.

The concrete frontier checker is deliberately separate from the exploratory
generator.  For each serialized layer and each input state it:
\begin{enumerate}[leftmargin=*,nosep]
  \item reconstructs the layer equations from the canonical manifest;
  \item asserts digit, leading-digit, carry, and input-state bounds in QF\_LIA;
  \item asks Z3 for a model, records the output projection, blocks that
        projection, and repeats until unsatisfiable; and
  \item compares the sorted projections, every input transition, live-cut
        variables, equation coverage, and content digests with the certificate.
\end{enumerate}
Z3's integer sort denotes mathematical integers, and its documented LIA
machinery uses cuts and branch-and-bound \cite{z3Arithmetic}.  The checker
rejects \texttt{unknown}; it never treats solver termination as
unsatisfiability.  Replaying the \Ccore{} certificate took
\CoreReplaySeconds{} seconds on the recorded development machine.

\begin{proposition}[computational obstruction]\label{prop:concrete}
Under the correctness of the manifest parser, the QF\_LIA encoding, and Z3's
reported results, \Ccandidate{} and \Ccore{} are unsatisfiable.
\end{proposition}

The qualifier in Proposition~\ref{prop:concrete} is material.  The Z3 runs do
not emit a proof trace checked by a small independent proof checker, and the
concrete frontier closure obligations are not discharged by the Lean kernel.

\section{Deletion-Minimal Core}

Let $e_\star$ be the equation \RemovedCoreEquation, and set
$\Ccore=\Ccandidate\setminus\{e_\star\}$.  The standalone frontier replay
establishes computational unsatisfiability of \Ccore.  A complete audit then
solves every system $\Ccandidate\setminus\{e\}$ for
$e\in\Ccandidate$.  Of \DeletionAuditCount{} outcomes,
\SatDeletionCount{} have complete satisfying assignments and
\UnsatDeletionCount{} is unsatisfiable; the latter is the deletion of
$e_\star$.

For any $e\in\Ccore$, the stored assignment for
$\Ccandidate\setminus\{e\}$ also satisfies its subset
$\Ccore\setminus\{e\}$.  Direct integer replay checks all bounds and retained
column equalities for every such assignment.  Therefore every one-equation
deletion of \Ccore{} is satisfiable in the same computational certificate
model, proving deletion minimality relative to the frozen candidate.

One neighboring result has a stronger endpoint.  Deleting
\KnownSatRemovedEquation{} leaves \CoreEquationCount{} equations, and the
full assignment is embedded as transparent Lean data.  Ordinary kernel
reduction checks the theorem \ConcreteSatTheorem{} in
\LeanWitnessSeconds{} seconds on the recorded machine; no
\texttt{native\_decide} oracle is used.  This kernel check covers one SAT
neighbor, not the complete deletion audit.

\begin{table}[t]
\centering
\small
\begin{tabular}{@{}lll@{}}
\toprule
Object & Equations & Assurance endpoint \\
\midrule
Candidate & \CandidateEquationCount & exact QF\_LIA frontier replay \\
Core & \CoreEquationCount & exact replay plus all-deletion witnesses \\
Known SAT neighbor & \CoreEquationCount & Lean kernel, ordinary \texttt{decide} \\
Generic soundness theorem & --- & Lean kernel \\
\bottomrule
\end{tabular}
\caption{Claim boundaries.  The first two rows are not Lean theorems.}
\label{tab:assurance}
\end{table}

\section{Reproducibility and Trust}

The artifact included with this preprint follows the computer-aided-proof
principles of making the computational role explicit and supplying
surveyable code and supporting data \cite{lmsComputerAidedProofs}.  Its quick
entry point verifies frozen SHA-256 values; regenerates the equation and Lean
data; independently replays both empty frontiers and all promoted SAT
witnesses; runs the publication tests; performs a clean Lean build; compares
the Lean-exported and Python-built statements; and prints theorem axiom
dependencies.  The command is
\begin{center}
\texttt{./publication/base19-radial/reproduce.sh quick}.
\end{center}
The recorded clean-clone run took \CleanCloneSeconds{} seconds.  The
\texttt{full} mode additionally regenerates the complete deletion audit and
all promoted certificates before byte comparison.  Certificate generation
alone took \CertificateGenerationSeconds{} seconds; the resumed deletion
audit took \DeletionAuditSeconds{} seconds.  Timings are observations, not
mathematical evidence.

The main trust paths are intentionally visible:
\[
\begin{array}{c}
\text{topology JSON}\longrightarrow\text{independent equation manifest}
\longrightarrow\text{layered QF\_LIA replay},\\[2mm]
\text{topology JSON}\longrightarrow\text{Lean equation builder}
\longrightarrow\text{canonical manifest comparison},\\[2mm]
\text{Lean definitions}\longrightarrow\text{generic theorem and one SAT witness}
\longrightarrow\text{Lean kernel}.
\end{array}
\]
The first path still trusts Z3, Python, the operating system, and the
arithmetic encoding.  The second catches translation drift but cannot turn a
solver status into a theorem.  The third has a smaller logical endpoint but
does not include the concrete UNSAT instance.

\section{Related Work}

Seal's construction and the OEIS entry provide historical reverse-and-add
context \cite{brownSeal,oeisA060382}.  Our radial frontier is an instance of
the familiar principle that bounded separators make structured CSPs
tractable \cite{Freuder1990,DechterPearl1989}; we make no claim that the
frontier method itself is new.

For solver-result certification, DRAT-trim checks expressive clausal proofs
\cite{WetzlerHeuleHunt2014}; LRAT adds hints suitable for simple certified
checkers \cite{CruzFilipeEtAl2017}; and FRAT supports communication from a
solver to a later elaborator \cite{BaekCarneiroHeule2022}.  A future Boolean
encoding could use a proof-producing solver such as CaDiCaL
\cite{cadical301}, followed by one of these checking paths.  We did not
promote such a translation or trace because it is not required by the direct
radial argument and was not validated to the standard of this artifact.

The literature search is documented in \texttt{paper/literature-matrix.md}.
It did not establish priority for the finite result, so this paper makes no
literature-wide novelty claim.

\section{Limitations}

The following boundaries are part of the result.
\begin{itemize}[leftmargin=*]
  \item The concrete \CandidateEquationCount-equation and
        \CoreEquationCount-equation unsatisfiability claims are not
        kernel-checked and have no independently checked LRAT, DRAT, or FRAT
        trace.
  \item Deletion minimality is computationally certified.  Only the one
        neighbor named above is separately checked in Lean.
  \item This is not a minimum-cardinality claim.  Smaller unsatisfiable
        subsets may exist.
  \item There is no all-base result and no proof that any particular starting
        integer is Lychrel.  The theorem concerns one fixed base-19 topology
        and a subset of its column equations.
  \item The one-deletion audit identifies a single unsatisfiable deletion of
        \Ccandidate, but this does not justify an unqualified uniqueness claim
        outside those immediate subsets.
  \item Runtimes and memory use are platform-dependent and are not part of
        the mathematical statement.
\end{itemize}

The most valuable next strengthening would be a compact concrete proof object
checked by a substantially smaller independent verifier, or a feasible Lean
proof of the concrete frontier closure.  Until then, Proposition
\ref{prop:concrete} should be cited with its computational assurance boundary.

\begingroup
\small
\bibliographystyle{plain}
\bibliography{references}
\endgroup

\end{document}
